\section{From 1D Signals to 3D Geometry: Evolution of Representation Learning}
\label{sec:signal-progression}

\subsection{1D Signals: Audio Processing and Sequential Data}

Audio waveforms were the first domain where the interplay between
sampling theory, transform design, and perceptual coding was worked out
rigorously, and the same ideas recur in 2D and 3D.
The Nyquist-Shannon sampling theorem~\cite{shannon1949noise} establishes
the foundational link between continuous and discrete representations:
a bandlimited signal $x(t)$ with maximum frequency $f_{\max}$ is perfectly
reconstructable from its samples $x[n]$ provided the sampling rate exceeds
$2f_{\max}$.
Without this result, the entire chain from ADC to digital storage to playback
is undefined.

Traditional audio codecs (MP3, AAC) exploit the mismatch between signal
information content and the resolution of human auditory perception.
The Modified Discrete Cosine Transform (MDCT)~\cite{princen1987mdct} decomposes overlapping frames
into frequency subbands; coefficients in less perceptually sensitive bands
are quantised more coarsely; the result is entropy-coded.
Designing the perceptual weighting functions required decades of
psychoacoustic research.

Deep learning replaced hand-crafted transforms with learned ones.
WaveNet~\cite{oord2016wavenet} modelled the raw audio waveform
autoregressively using dilated causal convolutions with exponentially growing
receptive fields, generating high-quality speech.
HiFi-GAN~\cite{kong2020hifigan} applied a generative adversarial network (GAN)
trained directly on waveforms to achieve real-time high-fidelity synthesis.
The core operation in all these systems is the 1D convolution:
a kernel $\mathbf{w} \in \R^k$ slides over the sequence
$x[n]$ to produce:
\begin{equation}
  y[n] = \sum_{i=0}^{k-1} w[i] \cdot x[n - i].
\end{equation}
Dilation by factor $d$ skips samples:
$y[n] = \sum_{i=0}^{k-1} w[i] \cdot x[n - d\cdot i]$,
exponentially expanding the receptive field without proportionally increasing
the parameter count.

\subsection{2D Signals: Image Understanding and Convolutional Networks}

Images extend 1D signals to a 2D regular grid $x[m, n]$ of pixel intensities.
The Discrete Cosine Transform (DCT), applied blockwise, underpins JPEG
compression; the H.264/H.265 family adds inter-frame motion compensation for
video.
Both systems are built on expert-crafted transforms: the DCT was chosen
for its energy-compaction properties; motion estimation is heuristic.

The revolution in image understanding came from \emph{learning} the transforms
directly from data.
A 2D convolution slides a learnable kernel
$\mathbf{W} \in \R^{k \times k \times C_{in} \times C_{out}}$
over the input to produce feature maps:
\begin{equation}
  y_{c'}[m, n] = \sum_{c=1}^{C_{in}} \sum_{i=-\lfloor k/2\rfloor}^{\lfloor k/2\rfloor}
    \sum_{j=-\lfloor k/2\rfloor}^{\lfloor k/2\rfloor}
    W_{c,c'}[i, j] \cdot x_c[m+i,\, n+j].
\end{equation}
The key structural property is \emph{weight sharing}: the same kernel is
applied at every spatial location, reducing parameter count from $O(N^2)$
to $O(k^2 C_{in} C_{out})$ relative to a fully connected layer.
A direct consequence is \emph{translation equivariance}: shifting the input
shifts the output by the same amount, so the network does not need to
re-learn a detector at every position.
Figure~\ref{fig:conv2d} illustrates the $3\times 3$ convolution operation.

\begin{figure}[t]
\centering
\begin{tikzpicture}[>=Stealth]
  \def\cs{0.56}  % cell size
  % ---- 5x5 input grid ----
  \foreach \i in {0,...,4}
    \foreach \j in {0,...,4} {
      \pgfmathsetmacro\clr{int(18+mod(\i*7+\j*11,20))}
      \fill[blue!\clr!white] ({\j*\cs},{-\i*\cs}) rectangle ++({\cs},{\cs});
      \draw[gray!50, very thin] ({\j*\cs},{-\i*\cs}) rectangle ++({\cs},{\cs});
    }
  % Receptive field: rows i=1..3, cols j=1..3 — fully inside the grid
  % x: [cs, 4*cs] = [0.56, 2.24];  grid x: [0, 5*cs=2.80]  ✓
  % y: [-3*cs, 0] = [-1.68, 0];    grid y: [-4*cs=-2.24, cs=0.56]  ✓
  \draw[red!80, very thick, rounded corners=1pt]
    ({\cs},{-3*\cs}) rectangle ++({3*\cs},{3*\cs});
  \node[font=\small, below] at ({2.5*\cs},{-5*\cs-0.06}) {Input $\mathbf{x}$};

  % Convolution symbol — at the vertical centre of the 5x5 grid
  \node[font=\large] at ({5*\cs+0.45},{-1.5*\cs}) {$*$};

  % ---- 3x3 kernel (yshift=-\cs centres it on the 5x5 input) ----
  \begin{scope}[xshift={5*\cs cm + 0.85cm}, yshift={-\cs cm}]
    \foreach \i in {0,...,2}
      \foreach \j in {0,...,2} {
        \fill[orange!28] ({\j*\cs},{-\i*\cs}) rectangle ++({\cs},{\cs});
        \draw[orange!60] ({\j*\cs},{-\i*\cs}) rectangle ++({\cs},{\cs});
        \node[font=\tiny] at ({\j*\cs+0.5*\cs},{-\i*\cs+0.5*\cs}) {$w_{\i\j}$};
      }
    \node[font=\small, below] at ({1.5*\cs},{-4*\cs-0.06}) {Kernel $\mathbf{w}$};
  \end{scope}

  % Equals sign
  \node[font=\large] at ({5*\cs+0.85+3*\cs+0.45},{-1.5*\cs}) {$=$};

  % ---- 3x3 output grid (yshift=-\cs centres it on the 5x5 input) ----
  \begin{scope}[xshift={5*\cs cm + 0.85cm + 3*\cs cm + 0.85cm}, yshift={-\cs cm}]
    \foreach \i in {0,...,2}
      \foreach \j in {0,...,2} {
        \fill[green!18] ({\j*\cs},{-\i*\cs}) rectangle ++({\cs},{\cs});
        \draw[green!55!black] ({\j*\cs},{-\i*\cs}) rectangle ++({\cs},{\cs});
      }
    % Highlighted output cell corresponding to the RF (row 1, col 1)
    \fill[green!55] ({\cs},{-\cs}) rectangle ++({\cs},{\cs});
    \node[font=\small, below] at ({1.5*\cs},{-4*\cs-0.06}) {Output $\mathbf{y}$};
  \end{scope}
\end{tikzpicture}
\caption{%
  A $3{\times}3$ convolution on a $5{\times}5$ input feature map.
  The kernel (orange) slides over the input, computing a dot product at each
  position to produce one output value (green).
  The highlighted receptive field (red border) shows the input region that
  contributes to one output position.
}
\label{fig:conv2d}
\end{figure}

Stacking convolutional layers with interleaved non-linearities and spatial
downsampling (max pooling, strided convolution) builds a hierarchy of features
whose receptive field grows with depth.
AlexNet~\cite{krizhevsky2012alexnet} was the first to demonstrate on ImageNet that this learned
hierarchy dramatically outperforms hand-crafted feature pipelines; VGG,
GoogLeNet, and ResNet~\cite{he2016deep} subsequently showed that depth and
efficient skip-connection design were the main levers for further gains.

The key architectural innovation that made very deep networks trainable was
the residual block.
Instead of learning a mapping $\mathbf{H}(\mathbf{x})$ from input to output,
a residual block learns the residual $\mathbf{F}(\mathbf{x}) = \mathbf{H}(\mathbf{x}) - \mathbf{x}$
and adds the identity:
\begin{equation}
  \mathbf{H}(\mathbf{x}) = \mathbf{F}(\mathbf{x}) + \mathbf{x}.
\end{equation}
This reformulation makes it easy for the network to learn the identity mapping
(by setting $\mathbf{F} = \mathbf{0}$), which ensures that adding more layers
cannot hurt performance.
Residual connections also provide a direct gradient path from the output to
early layers, alleviating the vanishing gradient problem.

For dense prediction tasks (segmentation, depth estimation, image restoration),
the output resolution must match the input.
U-Net~\cite{ronneberger2015unet} introduced the encoder-decoder with skip
connections: an encoder path progressively downsamples with strided
convolutions (or pooling), building a compact high-level representation;
a symmetric decoder path progressively upsamples; at each resolution level,
skip connections concatenate encoder features to decoder inputs, preserving
fine-scale spatial detail.
\begin{equation}
  \mathbf{h}^{(dec)}_\ell = f_\ell^{dec}\!\left(
    [\,\mathbf{h}^{(enc)}_\ell\, ;\, \text{upsample}(\mathbf{h}^{(dec)}_{\ell+1})\,]\right),
\end{equation}
where $[\cdot\,;\,\cdot]$ denotes channel concatenation.

Learned compression extended the same encoder-decoder idea to the
rate-distortion problem.
Ballé et al.~\cite{balle2018variational} replaced the hand-crafted DCT
with a learned CNN encoder-decoder and introduced the \emph{hyperprior}:
a second entropy model that codes the scale parameters of the main latent
distribution, capturing spatial correlations in the quantised latent.
Minnen et al.~\cite{minnen2018joint} added an autoregressive prior that
conditions each latent element on previously decoded context, further
improving entropy coding efficiency.
These systems now surpass JPEG-2000 on standard benchmarks and provided
the blueprint for learned point cloud codecs (Section~\ref{sec:compression}).

\subsection{3D Data: Geometric Representations and the Sparsity Challenge}

Extending image-like deep learning to 3D confronts an immediate structural
problem: the spatial data is \emph{sparse and irregular}.
A $512 \times 512$ image contains exactly $262\,144$ pixels in a regular grid.
The same scene at the same spatial resolution in 3D would occupy a
$512^3 \approx 134\,\text{M}$ voxel grid, of which fewer than 1\% are
occupied by surface points.
Applying standard 3D convolutions at this resolution requires $\sim$512 GB
of memory for feature maps at a modest 64 channels --- clearly impractical.

Processing only the occupied subset requires representations that make
sparsity explicit rather than storing an empty grid.
The three principal approaches to this are surveyed in
Section~\ref{sec:3d-representations}.


% ============================================================
